\title{Homework 2}
\date{{\em See Canvas for due date and time}}

% \paragraph{Submission instruction.}
% You need to submit the following file(s) to Canvas:
% \begin{itemize}
%     \item A pdf file named hw2.pdf containing your answers for problems 2 and 3.
    
%     \item A ipynb file named hw2.ipynb.
%     This file is modified from the original hw2.ipynb to include your answers for problem 1. 
% \end{itemize}

Please be aware of our {\em AI policy} per the course syllabus.

\section{Implementing MCTS}
{[6 points]}  Read and modify the provided hw2\_mcts.ipynb file to implement the following methods therein.
It is recommended to use lecture2.pdf on Canvas or an equivalent as a reference for your implementation.
\begin{enumerate}[label=(\alph*)]
\item {[2 points]} \texttt{tree\_policy} in class \texttt{MCTS} when transition $(s,a,s')$ is already in the tree.
\item {[1 point]} \texttt{best\_action} in class \texttt{MCTS}.
\item {[1 point]} \texttt{search} in class \texttt{MCTS}.
\item {[1 point]} \texttt{select\_action} in class \texttt{UCT}.
\end{enumerate}
The original hw2\_mcts.ipynb file provides a test at the end, which only serves as a reference for grading.
After implementing all methods above, run the test and answer the following question:
\begin{enumerate}[label=(\alph*)]
\setcounter{enumi}{5}
\item {[1 point]} Can you briefly summarize your test result and interpret it?
\end{enumerate}

\section{Analysis of MAB Algorithms}
{[4 + 4 points]} Read lecture3a.pdf on Canvas carefully.
Answer the questions tagged as 2a-2d in lecture3a.pdf, 2 points each.
The questions are \ul{\em underlined and in italics}.

In particular, 2c and 2d are more challenging, they will be considered as bonus questions (4 points in total) for this homework.


\section{Policy Optimization}
{[3 points]} Read lecture3c.pdf on Canvas and justify the derivations tagged as {\bf 3a-3c} therein, 1 points each.

\section{Implmenting PPO}
{[7 points]} Read hw2\_ppo.ipynb that implements a simplified version of PPO.
Fill in your implementation of the unfinished methods therein.
Details can be found in hw2\_ppo.ipynb.



